\chapter{The Bracketed Number}

\section{The bound is the object}

The slip became a reading, and the reading was an interval. This chapter makes that interval the instrument's permanent
output --- not a temporary state of ignorance on the way to a point, but the object itself. The instrument owns the bracket
$[\ell, u]$ and does not own any point inside it. This is a reversal of the usual order of things, in which the value is
primary and the error bars are an apology attached afterward. Here the bracket is primary. The value is the fiction; the
interval is the fact.

State the reversal exactly, because it governs everything downstream. In the naive account a measurement is a number
$x$ with an uncertainty $\pm\varepsilon$, and the number is what you believe while the uncertainty is what you concede. The
instrument inverts the emphasis: what it holds is the pair of walls $\ell < u$, and the ``number'' $x$ is at most a
convenient midpoint with no independent standing. There is no $x$ in the construction that the walls approximate; there are
the walls, and the interval they bound, and nothing between them the instrument claims to know. The bracket is not $x \pm
\varepsilon$ with $x$ real and $\varepsilon$ regrettable. It is $[\ell, u]$ with the interval real and the midpoint a mere
label of convenience.

The reason the bound and not the point is the object is the reason the reals were built as approached-never-completed. A
point would be a completed real, plucked whole from the continuum --- and the continuum is exactly the object the finite
instrument declines to assume. What the instrument can actually produce is a lower rational and an upper rational, each a
finite condition, with a certificate that the true value lies between. That certificate is the reading. Ask the instrument
for more --- for the point --- and it has nothing more to give, because the point was never among its finite conditions. The
bracket is not the instrument stopping short of the point; it is the instrument reporting the whole of what its finite
conditions contain. There is no further fact it is withholding.

This has a consequence for how every later magnitude must be read, and it is worth fixing now. When the book reaches the
coupling, the coupling will be a bracket, and it will be tempting --- because the reader is trained to want points --- to
read the bracket as an approximation to a hidden true value the instrument almost reached. That reading will be wrong, and
wrong in the precise way this section forbids: the bracket is not an approximation to a point the instrument owns; it is the
object the instrument owns, and there is no point behind it in the instrument's possession. The world may have a value; the
instrument has an interval; and the honest report is the interval, full stop.

In the two verbs, the bracket is the funge the instrument completes and the point is the tange it declines. It funges the
value into the interval $[\ell, u]$ --- gathers it, honestly, between two certified walls --- and refuses the final tange
that would select a single real from the bag. The interval is the funge it can stand behind; the point is a selection its
finite conditions cannot support. So the bound is the object because the funge is the thing the instrument actually holds,
and the point is a tange it can neither perform nor honestly fake. What sets the width of that funge --- the floor below
which the walls cannot be driven any closer --- is the subject of the next section.

\section{The floor}

The bracket has a width, and the width is not arbitrary: below a certain scale the instrument cannot drive its two walls any
closer, because below that scale nothing registers. This section is about that floor --- the resolution beneath which the
loop's slip is indistinguishable from no slip, the cut that sets the bracket's width, and the tail below it that is the
instrument's own ignorance made explicit. The floor is the most physical object in the chapter, and on this face it has a
temperature.

Begin with why there is a floor at all. The instrument reads a residue by a loop, and a loop can resolve a difference only
if the difference exceeds the noise of the frame the loop is carried in. The construction takes its reading frame to be an
\emph{accelerated} one --- a uniformly accelerated observer, the Rindler frame~\cite{rindler2006} --- and an accelerated
frame is not quiet. It is warm. Unruh showed that a detector accelerated through the vacuum registers a thermal bath at
temperature
\[
  k_B T \;=\; \frac{\hbar\,a}{2\pi c},
\]
proportional to the acceleration $a$~\cite{unruh1976}. That temperature is the floor. A slip smaller than the thermal jitter
of the frame cannot be told from the jitter; the loop returns it inside the noise, and the instrument, honestly, cannot
resolve it. The resolution floor is not a limitation of a particular bench; it is the temperature of the frame the reading is
taken in, and it is set by the acceleration the frame carries.

The floor must be approached the right way, and the construction is emphatic about this because the wrong way is the crank's
way. A hard cut --- a sharp step that keeps everything above the floor and discards everything below --- \emph{rings}. Cut a
signal with a sharp edge and it develops oscillations near the edge, the Gibbs phenomenon, a spurious structure that is an
artifact of the cut and not a feature of the signal. An instrument that reads its floor with a hard step will manufacture
exactly such artifacts and mistake them for readings. The honest response is to \emph{taper} --- to approach the floor
smoothly, apodizing the cut, so that the transition from resolved to unresolved is gradual and rings nowhere. The accelerated
frame supplies the taper for free: the thermal floor is not a sharp wall but a smooth Boltzmann roll-off, $\exp(-E / k_B T)$,
which is exactly the apodization a hard cut lacks. The instrument tapers to its floor because its floor is thermal, and a
thermal floor is smooth by nature. Taper to the floor; never step to it.

Below the floor lies a tail, and the tail is the instrument's ignorance written down. The sub-floor residue does not vanish
--- it is there, smaller than the resolution, contributing something the instrument cannot resolve. The honest instrument
does not pretend the tail is zero; it marks it as unresolved and carries its existence as the reason the bracket has the
width it does. The width $u - \ell$ is, precisely, the size of what lives below the floor: everything the instrument cannot
tell from noise. So the bracket's width is not a failure of engineering to be reduced by a better bench. It is the published
size of the sub-floor tail, the instrument's own statement of the ignorance it cannot resolve away. A narrower bracket would
be a claim to have resolved the tail, and the tail, by construction, is what cannot be resolved.

In the two verbs, the floor is where tange runs out. Above it, the instrument can tange --- select one candidate slip from
another, tell them apart, drive the walls together. At the floor the thermal jitter erases the distinction: below it no
characteristic separates one candidate from the next, so no tange is available, and the instrument funges everything below
the floor into a single unresolved tail. The floor is the exact scale at which selection fails and only gathering remains.
That is why the bracket cannot be narrowed past it: to narrow it would be to tange below the floor, to select where nothing
distinguishes, and the instrument refuses that selection because the frame's own temperature forbids it. The last section
takes this floored bracket and shows that its honesty is repeatability --- the same interval, returned every time, because
the floor is universal.

\section{Repeatability is the honesty}

The bracket has a floor, and the floor has a temperature. This closing section draws out why a reading built this way is
trustworthy at all, and the answer is repeatability: a reading is the bracket you can run again and get back the same. But
the deep content of the section is \emph{why} the bracket repeats --- and the reason is that the floor is not a property of a
particular apparatus but a universal invariant, the same in every frame that reads at the same acceleration. Repeatability is
not luck; it is the universality of the floor.

Take repeatability first at face value. Run the loop, read the slip, report $[\ell, u]$; run it again and report the same
$[\ell, u]$. A reading that changed on each run would be no reading. What makes the interval stable is that both of its
determinants are stable: the residue it brackets is the invariant curvature, unchanged by relabeling, and the floor that sets
its width is the thermal temperature of the reading frame, fixed by the acceleration. Neither the value being read nor the
resolution it is read at depends on the arbitrary labels, so the interval comes back the same. Repeatability is the joint
invariance of the residue and the floor, seen through the loop.

Now the deeper claim, which is the reason four gauges can ever agree. The Unruh floor is \emph{universal}: any detector
accelerated at $a$ through the vacuum reads the same temperature $k_B T = \hbar a / 2\pi c$, regardless of what the detector
is made of or how it is built~\cite{unruh1976}. The floor does not know which apparatus reads it. This is the cross-setup
invariance the whole instrument is built to earn --- the fact that the resolution at which the bracket closes is a property
of the frame, not the bench, so two entirely different apparatuses, reading at the same acceleration, floor their brackets at
the same width. A quantity floored at a universal temperature is a quantity two independent instruments must bracket alike.
That universality is what will let the physical gauge and its three siblings, sharing no vocabulary, close on one interval:
they read the same residue at the same universal floor, so they cannot help but agree.

This is also where the honesty of the whole method becomes visible as a physical, not merely a moral, property. An instrument
that reported a point could tune that point to whatever it wished, and no repetition would catch it, because a point can be
re-announced unchanged forever regardless of whether it is true. An instrument that reports a floored bracket cannot cheat
this way: the floor is set by the frame's temperature, which the instrument does not control, so the width of its interval is
imposed on it from outside its own choices. It cannot narrow the bracket without lowering the acceleration, and it cannot
lower the acceleration without changing the frame it reads in. The honesty is enforced by physics: the bracket is as wide as
the floor makes it, the floor is as warm as the acceleration makes it, and none of it is the instrument's to fake.

In the two verbs, repeatability is the funge that returns the same bag every time, because the floor that sets the bag is the
world's and not ours. The instrument funges the reading into $[\ell, u]$; the width of that funge is fixed by a covariant,
universal temperature no tange of the labels can alter; so the funge is stable, and rerunning it returns the identical
interval. The point a crank would tange out of the bag is the instrument's to invent and therefore to fake; the interval the
floor funges is the world's to impose and therefore to trust. So Part II closes with the instrument's output fully
characterized: a bracketed, dimensional, repeatable reading of an invariant residue, floored at a universal temperature, held
open by a tail it honestly cannot resolve. What remains is to ask what that single bracketed residue \emph{is} --- and the
answer, next, is that it is not one thing but four faces of one thing, depending on where the corridor is read.
